\chapter{2014年安徽卷（文）}

\section{单选题}

\begin{problem}
设 1 是虚单位基，复数 \( \i^{2} + \frac{2\i}{1+1} = \)（\quad）

\choices
    {-i}
    {i}
    {\(-1\)}
    {1}
\end{problem}

\begin{problem}
命题“\(\forall x \in \R, |x| + x^2 \geqslant 0\)”的否定是（\quad）

\choices
    {\(\forall x\in \R,|x| + x^2 < 0\)}
    {\(\forall x\in \R,|x| + x^2\leqslant 0\)}
    {\(\exists x_0\in \R,|x_0| + x_0^2 < 0\)}
    {\(\exists x_0\in \R,|x_0| + x_0^2\geqslant 0\)}
\end{problem}

\begin{problem}
抛物线 \( y = \frac{1}{4} x^2 \) 的准线方程是（\quad）

\choices
    {\(y = -1\)}
    {\(y = -2\)}
    {\(x = -1\)}
    {\(x = -2\)}
\end{problem}

\begin{problem}
如图所示，程序框图 (算法流程图) 的输出结果是（\quad）

\choices
    {34}
    {55}
    {78}
    {89}
\end{problem}

\begin{problem}
设 \(a = \log_6 T, b = 2^{1.1}, c = 0.8^{0.1}\)，则（\quad）

\choices
    {\(b < a < c\)}
    {\(c < a < b\)}
    {\(c < b < a\)}
    {\(a < c < b\)}
\end{problem}

\begin{problem}
过点 \(P(- \sqrt{3}, -1)\) 的直线 \(l\) 与圆 \(x^{2} + y^{2} = 1\) 有公共点，则直线 \(l\) 的倾斜角的取值范围是（\quad）

\choices
    {\(\left(0, \frac{\pi}{6}\right]\)}
    {\(\left(0, \frac{\pi}{3}\right]\)}
    {\(\left[0, \frac{\pi}{6}\right]\)}
    {\(\left[0, \frac{\pi}{3}\right]\)}
\end{problem}

\begin{problem}
若将函数 \( f(x) = \sin 2x + \cos 2x \) 的图象向右平移 \( \varphi \) 个单位，所得图象关于 \( y \) 轴对称，则 \( \varphi \) 的最小正值是（\quad）

\choices
    {\(\frac{\pi}{8}\)}
    {\(\frac{\pi}{4}\)}
    {\(\frac{3\pi}{8}\)}
    {\(\frac{3\pi}{4}\)}
\end{problem}

\begin{problem}
一个多面体的三视图如图所示，则该多面体的体积为（\quad）

正(主)视图

侧(左)视图

俯视图（\quad）

\choices
    {\( \frac{23}{3} \)}
    {\( \frac{47}{6} \)}
    {6}
    {7}
\end{problem}

\begin{problem}
若函数 \( f(x) = |x + 1| + |2x + a| \) 的最小值为3，则实数 \( a \) 的值为（\quad）

\choices
    {5 或 8}
    {\(-1\) 或 5}
    {\(-1\) 或 \(-4\)}
    {\(-4\) 或 8}
\end{problem}

\begin{problem}
设 \(a, b\) 为非零向量，\(|b| = 2 |a|\)，两组向量 \(x_1, x_2, x_3, x_4\) 和 \(|y_1, y_2, y_3, y_4\) 均由 2 个 \(a\) 和 2 个 \(b\) 排列而成，若 \(x_1 \cdot y_1 + x_2 \cdot y_2 + x_3 \cdot y_3 + x_4 \cdot y_4\) 所有可能取值中的最小值为 \(4 |a|^2\)，则 \(a\) 与 \(b\) 的夹角为（\quad）

\choices
    {\(\frac{2}{3}\pi\)}
    {\(\frac{\pi}{3}\)}
    {\(\frac{\pi}{6}\)}
    {0}
\end{problem}

\section{填空题}

\begin{problem}
\(\left(\frac{16}{81}\right)^{-1} + \log_5\frac{4}{9} + \log_5\frac{4}{5} =\)  \fillinblank{}.
\end{problem}

\begin{problem}
如图，在等腰直角三角形 \(ABC\) 中，斜边 \(BC = 2\sqrt{2}\)，过点 \(A\) 作 \(BC\) 的垂线. 垂足为 \(A_{1}\); 过点 \(A_{1}\) 作 \(AC\) 的垂线，垂足为 \(A_{2}\); 过点 \(A_{2}\) 作 \(A_{1}C\) 的垂线. 垂足为 \(A_{3}\); ……，依此类推. 设 \(BA = a_{1}\)，\(AA_{1} = a_{2}\)，\(AA_{1}A_{2} = a_{3}\)，\(\cdots\)，\(A_{n}A_{n} = a_{7}\)，则 \(a_{7} =\)  \fillinblank{}.
\end{problem}

\begin{problem}
不等式组 \(\left\{ \begin{array}{l} x + y - 2 \geqslant 0, \\ x + 2y - 4 \leqslant 0, \end{array} \right.\) 表示的平面区域的面积为 \( (x + 3y - 2 \geqslant 0 \)
\end{problem}

\begin{problem}
若函数 \( f(x) \) (\( x \in \R \)) 是周期为 4 的奇函数，且在 \([0, 2]\) 上的解析式为 \( f(x) = \begin{cases} x(1 - x), & 0 \leqslant x \leqslant 1, \\ \sin \pi x, & 1 < x \leqslant 2. \end{cases} \) 则 \( f\left(\frac{29}{4}\right) + f\left(\frac{41}{6}\right) = \)  \fillinblank{}.
\end{problem}

\begin{problem}
若直线 \(l\) 与曲线 \(C\) 满足下列两个条件: (i) 直线 \(l\) 在点 \(P(x_0, y_0)\) 处与曲线 \(C\) 相切；(ii) 曲线 \(C\) 在点 \(P\) 附近位于直线 \(l\) 的两侧，则称直线 \(l\) 在点 \(P\) 处“切过”曲线 \(C\). 下列命题正确的是 \fillinblank{}. (写出所有正确命题的编号) \circled{1} 直线 \(l: y = 0\) 在点 \(P(0,0)\) 处“切过”曲线 \(C: y = x^2\); \circled{2} 直线 \(x: x = -1\) 在点 \(P(-1,0)\) 处“切过”曲线 \(C: y = (x + 1)^2\); \circled{3} 直线 \(l: y = x\) 在点 \(P(0,0)\) 处“切过”曲线 \(C: y = \sin x\); \circled{4} 直线 \(l: y = x\) 在点 \(P(0,0)\) 处“切过”曲线 \(C: y = \tan x\); \circled{5} 直线 \(l: y = x - 1\) 在点 \(P(1,0)\) 处“切过”曲线 \(C: y = \ln x\) \fillinblank{}.
\end{problem}

\section{解答题}

\begin{problem}
设 \(\triangle ABC\) 的内角 \(A, B, C\) 所对边的长分别是 \(a, b, c\)，且 \(b = 3, c = 1\)，\(\triangle ABC\) 的面积为 \(\sqrt{2}\)，求 \(\cos A\) 与 \(a\) 的值.
\end{problem}

\begin{problem}
某高校共有学生 15000 人，其中男生 10500 人，女生 4500 人，为调查该校学生每周平均体育运动时间的情况. 采用分层抽样的方法，收集 300 位学生每周平均体育运动时间的样本数据 (单位: 小时).

(1) 应收集多少位女生样本数据\(\perp\) (2) 根据这 300 个样本数据，得到学生每周平均体育运动时间的频率分布直方图（如图所示），其中样本数据分组区间为：[0,2]，(2,4]，(4,6]，(6,8]，(8,10)，(10,12]. 估计该校学生每周平均体育运动时间超过 4 个小时的概率. (3) 在样本数据中，有 60 位女生的每周平均体育运动时间超过 4 个小时. 请完成每周平均体育运动时间与性别的列联表，并判断是否有 95\% 的把握认为“该校学生的每周平均体育运动时间与性别有关”. 附： \(K^2 = \frac{n(ad - bc)^2}{(a + b)(c + d) + (a + c)(b + d)}\) \( P(K^{2} \geqslant k_{0}) \) 0.10 0.05 0.010 0.005 \( k_{0} \) 2.706 3.841 6.635 7.879
\end{problem}

\begin{problem}
数列 \(\{a_{n}\}\) 满足 \(a_1 = 1, na_{n+1} = (n + 1)a_n + n(n + 1), n \in \mathbf{N}^*\). (1) 证明: 数列 \( \left\{\frac{a_{n}}{a_{n}}\right\} \) 是等差数列; (2) 设 \( b_n = 3^n \)、\( \sqrt[n]{a_n} \)，求数列 \( \{b_n\} \) 的前 \( n \) 项和 \( S_n \).
\end{problem}

\begin{problem}
设函数 \( f(x) = 1 + (1 + a)x - x^2 - x^3 \)，其中 \( a > 0 \). (1) 讨论 \( f(x) \) 在其定义域上的单调性; (2) 当 \(x \in [0,1]\) 时，求 \(f(x)\) 取得最大值和最小值时的 \(x\) 的值.
\end{problem}

\begin{problem}
设 \(F_{1}, F_{2}\) 分别是椭圆 \(E: \frac{x^{2}}{a^{2}} + \frac{y^{2}}{b^{2}} = 1 (a > b > 0)\) 的左、右焦点，过点 \(F_{1}\) 的直线交椭圆 \(E\) 于 \(A, B\) 两点，\(|AF_{1}| = 3|BF_{1}|\). (1) 若 \( \left|AB\right|=4,\triangle ABF_{2} \) 的周长为 16，求 \( \left|AF_{2}\right| \) ; (2) 若 \(\cos \angle AF_2B = \frac{3}{5}\)，求椭圆 \(E\) 的离心率.
\end{problem}

\begin{problem}
如图，四棱锥 \(P - ABCD\) 的底面是边长为 8 的正方形，四条侧棱长均为 \(2\sqrt{17}\)，点 \(G, E, F, H\) 分别是棱 \(PB, AB, CD, PC\) 上共面的四点，平面 \(GEFH \perp\) 平面 \(ABCD, BC \not//\) 平面 \(GEFH\). (1) 证明: \(GH // EF\); (2) 若 \(EB = 2\)，求四边形 \(GEFH\) 的面积.
\end{problem}

